%% Generated by manuscript/build/build.py -- edit the source, not this file.
%% PLOS LaTeX template (journals.plos.org/ploscompbiol/s/latex), adapted for XeTeX/Tectonic.
\documentclass[10pt,letterpaper]{article}
\usepackage[top=0.85in,left=1in,right=1in,footskip=0.75in]{geometry}
\usepackage{amsmath,amssymb}
\usepackage{changepage}
\usepackage{textcomp,marvosym}
\usepackage{cite}
\usepackage{nameref,hyperref}
\usepackage[right]{lineno}
\usepackage{iftex}
\usepackage[nopatch=eqnum]{microtype}
\ifpdftex\DisableLigatures[f]{encoding = *, family = * }\fi
\usepackage[table]{xcolor}
\definecolor{boxgray}{gray}{0.94}
\usepackage{array}
\usepackage{fvextra}
\usepackage{url}
\newcolumntype{+}{!{\vrule width 2pt}}
\newlength\savedwidth
\newcommand\thickcline[1]{%
  \noalign{\global\savedwidth\arrayrulewidth\global\arrayrulewidth 2pt}%
  \cline{#1}%
  \noalign{\vskip\arrayrulewidth}%
  \noalign{\global\arrayrulewidth\savedwidth}%
}
\newcommand\thickhline{\noalign{\global\savedwidth\arrayrulewidth\global\arrayrulewidth 2pt}%
\hline
\noalign{\global\arrayrulewidth\savedwidth}}
\raggedright
\setlength{\parindent}{0.5cm}
\textheight 8.75in
\usepackage[aboveskip=1pt,labelfont=bf,labelsep=period,justification=raggedright,singlelinecheck=off]{caption}
\renewcommand{\figurename}{Fig}
\bibliographystyle{plos2025}
\makeatletter
\renewcommand{\@biblabel}[1]{\quad#1.}
\makeatother
\usepackage{lastpage,fancyhdr,graphicx}
\pagestyle{fancy}
\fancyhf{}
\rfoot{\thepage/\pageref{LastPage}}
\renewcommand{\headrulewidth}{0pt}
\renewcommand{\footrule}{\hrule height 2pt \vspace{2mm}}
\lfoot{\today}
%% PLOS asks for figures NOT to be embedded in the submission file; set \embedfiguresfalse for submission.
\newif\ifembedfigures
\embedfiguresfalse
\begin{document}
\begin{flushleft}
{\Large\textbf{S3 Text. Extended methods and results}}
\end{flushleft}
This text holds the full technical description of BURROW, the benchmarks and the statistics, and the longer versions of several Results paragraphs that were shortened in the main text. Reference numbers in square brackets refer to the reference list of the main text. Sources: docs/METHOD.md, docs/TRAINSET.md, docs/EVALUATION.md, docs/DIMER\_BENCHMARK.md, docs/AF\_MODE.md, docs/PERF.md, docs/PREREG.md and REPORT.md in the project repository.

\section*{1. Extended methods}

\subsubsection*{Input and candidate points}

BURROW reads PDB or mmCIF files with gemmi [36] and keeps protein heavy atoms (first model, first alternate location). Ligands, waters, ions, glycans and nucleic acids are removed, so predictions do not depend on HETATM records (tested). Geometry is computed in a canonical frame given by the principal axes. Grid points (1 \AA{} spacing) are kept if they are at least the van der Waals radius plus 1 \AA{} from every atom and within 6 \AA{} of the nearest atom. For each point, BURROW counts LIGSITE protein--solvent--protein events along 7 directions up to 8 \AA{} [1]. Points with a count of 2 or more are kept at full resolution (volume 1 \AA{}$^{3}$), and the rest on a 3 \AA{} sub-lattice. Metal candidates are placed at 1.95--2.35 \AA{} from a side-chain atom that can coordinate a metal. Coordinating atoms are His ND1/NE2, Cys SG, Asp or Glu carboxylate O, Met SD and backbone O. A candidate is kept if at least two coordinating atoms (at least one from a side chain) lie within 2.6 \AA{} of it and no other atom lies closer than 2.4 \AA{}. They are deduplicated at 1.5 \AA{}, with at most one per cluster of coordinating atoms.

\subsubsection*{Point features}

Rays in 30 fixed Fibonacci directions (60 in the accurate preset) are marched through two absorption grids: all heavy atoms, and backbone plus C$\beta$ atoms only. The occlusion radius is the van der Waals radius plus 0.5 \AA{}. For each grid, BURROW computes the blocked fraction at ray lengths of 6, 10 and 16 \AA{} (plus 24 \AA{} in the accurate preset), the mean depth of the first hit, the fractions of first hits by chemical class, and the escape fraction. Other features are the protein--solvent--protein count, the nearest-atom distance, a metal flag with the number and angular spread of coordinators within 2.8 \AA{}, atom counts at 4--12 \AA{}, counts of 13 atom classes and 7 residue groups at 4 \AA{} and their fractions at 8 and 12 \AA{}, the mean weighted contact number (sum of 1/d$^{2}$ over C$\alpha$ atoms) at 8 and 12 \AA{}, and protrusion (atoms within 10 \AA{}). This gives 84 features (86 in the accurate preset). The geometric value of a point, b, is the mean blocked fraction at 10 \AA{}, 16 \AA{} and the longest ray length (so 16 \AA{} counts twice in the default preset; 10, 16 and 24 \AA{} in the accurate preset).

\subsubsection*{Point scorer and training labels}

The point scorer is a LightGBM binary classifier [24] with 250 trees of 63 leaves, bagging and feature fractions of 0.8, and fixed seeds. For the \texttt{all} ranker, a point is positive if it lies within 4 \AA{} of a non-ion ligand atom or within 3 \AA{} of an ion. For the \texttt{small} ranker, it is positive if it lies within 4 \AA{} of a non-ion ligand atom. Training samples 30\% of positive and 3\% of negative points per structure. In the default preset, only points with b $\ge$ 0.7 and metal candidates are scored (about 45\% of points). Gradient-boosted trees were chosen because they are fast on tabular features and because the prompt warned that large deep-learning models overfit. Most hyperparameters (grid spacing, number and length of rays, thresholds, tree sizes) were set in early runs on the P2Rank training sets (CHEN11 and JOINED) and then chosen among about ten settings on dev under the speed budget of decision D5; ranker variants were selected on training-set cross-validation. Dev numbers are therefore somewhat optimistic; test sets were not used for any choice.

\subsubsection*{Pocket formation, ranking and suppression}

Candidate pockets are the union of a geometric pass (single-linkage components at 2.05 \AA{} of points with b $\ge$ 0.40, 0.55, 0.70 or 0.85) and a score pass (point score $\ge$ 0.15, 0.30 or 0.50). Components longer than 20 \AA{} are split by weighted k-means into pieces of about 12 \AA{}, and components longer than 10 \AA{} are also split into pieces of about 6 \AA{}. Pockets smaller than 10 \AA{}$^{3}$ are dropped (metal candidates excepted). The centre is the centroid of the points weighted by volume and geometric value. The pocket ranker is a LightGBM model (300 trees, 15 leaves) on 76 descriptors: aggregates and quantiles of point scores, volume, buriedness (including backbone-only), first-hit chemistry, residue groups, shape from principal components, number of residues and chains, and protein size. Its training uses cross-fitted point scores from 5 folds of 30\% sequence clusters. The ranker is trained with cross-entropy on the target clip(1 $-$ q/12, 0, 1), where q is the smallest DCC to a site whose ligand atoms lie within 4 \AA{} of the pocket centre. Suppression is greedy in order of ranker score. A pocket is removed if its centre is within 10 \AA{} of a kept pocket, its residue set has a Jaccard index of 0.75 or more with one, or its point overlap with one is 0.3 or more. At most 50 pockets are kept. In v2, pockets with score below $\tau$ = 0.03 are dropped once three have been kept. $\tau$ was chosen by cross-validation over clusters of the training set, under the condition that closest-pocket recall, DCA4 recall, one-to-one P1K and top-N must not decrease. Reported residues are those with a heavy atom within 4 \AA{} of a pocket point.

\subsubsection*{Presets, AlphaFold mode and multimers}

The default preset (30 rays, 3 \AA{} sub-lattice, partial scoring) was used for all pre-registered claims. The accurate preset (60 rays, 2 \AA{} sub-lattice, all points scored) is shipped because it was more than 1 point better on non-ion dev sites. The fast preset (1.5 \AA{} grid, 12 rays, 2 \AA{} sub-lattice, a point model of 100 trees with 31 leaves) was added after the test runs. It was selected on training-set cross-validation with a rule fixed in advance (at most 35 ms per structure and at most a 2-point loss) and confirmed once on dev.

\texttt{--af-mode} takes residue pLDDT from the B-factor column of the C$\alpha$ atom. It removes terminal runs with pLDDT \textless{} 50 and internal runs of 7 or more residues with pLDDT \textless{} 50 [21]. Remaining atoms get an opacity w = clip((pLDDT $-$ 40)/40, 0, 1) in both absorption grids, and atoms with w \textless{} 0.5 do not exclude candidate points. The standard models are used unchanged. An optional ranker (\texttt{--af-ranker}) adds pocket pLDDT descriptors and was trained on AFDB models of training proteins with transferred labels. A structure-only stand-in for inter-domain uncertainty, based on spectral bisection of the C$\alpha$ contact graph, gave no gain and is not used.

Multimers need no special mode. All chains enter the grids, features and ranker, and one pocket descriptor is the number of chains that line the pocket.

\subsubsection*{CPU and GPU backends}

The reference implementation (NumPy [37], SciPy [38]) is the specification. The fast backend uses Numba on the CPU [39] and Triton/PyTorch kernels on the GPU [40, 41], with ray marching in double precision and a tree-inference routine that does not depend on LightGBM. GPU runs use one process per GPU, with CPU workers for parsing and post-processing. Outputs (scores, centres, residues and points) were identical to the reference on 64 random dev structures for four models on CPU and GPU, and on 40 AlphaFold models in AlphaFold mode. A bug that passed one threshold to a GPU kernel in single precision was found after the v1 freeze. It changed outputs only for atoms at pLDDT exactly 70. It was fixed with a regression test, and all v1 AlphaFold-mode test predictions were generated on the CPU path.

\subsubsection*{Training and development data}

Training data came from BioLiP2 (snapshot 2025-11-12) [26]: X-ray or EM structures at 3.0 \AA{} or better, receptor chains of at least 50 residues, and BioLiP2's relevant ligands including metal ions. Before the training set was built, the agents froze and checksummed the union of all test sets. This union covered LIGYSIS, all polymer entities of COACH420 and HOLO4K, all CryptoBench splits, PocketMiner, the PocketWatch test set, the AlphaFold set and every dimer chain. A training chain was removed if an MMseqs2 search [42] found any hit to the test union with E \textless{} 10$^{-3}$ and identity of at least 30\% over at least 50\% coverage or at least 80 aligned residues, or if it shared a PDB ID or UniProt accession with a test entry. This removed 77.4\% of BioLiP2 chains. Sites were built in the LIGYSIS style [11]. For each UniProt accession, a representative chain was chosen, and ligands from up to 20 other entries were superposed onto it by local, sequence-anchored superposition (C$\alpha$ RMSD $\le$ 2 \AA{}). Ligands were then clustered by the relative intersection of their binding-residue fingerprints with average linkage, cut at 0.5. Clashing non-ion ligand atoms were removed. Before the audit fix, the training set had 14,073 sites, of which 37.7\% contained only ions. Dev data were 10\% of sequence clusters, built in the same way and also leakage-filtered against training. The independent audit found 7 assembly-partner chains that violated the rule, 92 evaluated inputs that violated it when sequences of observed residues were used (88 of them at 30--35\% identity), and one UniProt overlap. Of these, 8 were hard failures; the agents removed every entry that any tier of the audit flagged (186 entries, 2.6\%), which left 6,475 training and 579 dev chains. No full re-audit was run after this fix. The test verifier later found none of the 2,565 LIGYSIS UniProt accessions in training or dev data.

\subsubsection*{Benchmarks, metrics and baselines}

Table 2 of the main text lists the benchmarks. LIGYSIS was used with the authors' ligand-free representative chains and their redundancy filter, which leaves 6,721 of 6,882 sites for scoring [11]. Under top-N+2 recall, a site counts as found if the pocket whose centre is closest to the site centre has DCC $\le$ 12 \AA{} and rank $\le$ N+2, where N is the number of sites on the chain. C1-strict uses ligand atoms placed on the representative chain by superposition (6,387 of 6,721 sites). P1K is the precision of the 1,000 highest-scoring pockets pooled over the dataset. COACH420 and HOLO4K were scored with P2Rank's rules (DCA $\le$ 4 \AA{}, top-n) using the explicit relevant-ligand (mlig) versions [5]; our evaluator reproduced P2Rank's own evaluator to 0.1\%. The AlphaFold set maps LIGYSIS sites onto AFDB v6 models by UniProt numbering, after a check of residue types. The dimer set uses assembly 1 of entries of LIGYSIS proteins, restricted to assemblies defined by both authors and software, X-ray or EM at 3 \AA{} or better, and 2--8 protein chains. A site is at an interface if 3 or more residues on each of 2 or more chains contact it. Symmetric copies are grouped at a residue-set Jaccard index of 0.5 or more, and every method's predictions are collapsed by the same rule before scoring. scPDB [43] was not used as a test set because it has no stated license, one site per entry, and overlap with LIGYSIS. AF2BIND data were not used because they are residue-level with about 65 test chains, and HOTPocket because it is a catalogue of predictions, not a benchmark.

Comparators were run on identical, cleaned inputs: P2Rank 2.5.1 (default, and its AlphaFold configuration on predicted structures) [5] and fpocket 4.2.3 rescored by PRANK [3, 4]. On LIGYSIS we used PRANK 2.4.2, as in the original paper. Both versions gave identical LIGYSIS numbers. Published per-chain predictions of the deep-learning methods on LIGYSIS were scored with our harness and are descriptive, because they have no paired CI. Several deep-learning tools were also wrapped and run on some sets, but their training data overlap with the legacy benchmarks, so they were descriptive only (S1 Text).

\subsubsection*{Statistics}

The bootstrap resampled sequence clusters (30\% identity, MMseqs2) 1,000 times with a fixed seed. Multi-chain structures mapped to the union of their chains' clusters. All differences are paired, failures count as misses, and P1K is recomputed in each resample. Family 1 had three requirements, each for every seed and against both the published and our re-run fpocket+PRANK: for C1, a difference of at least +2 points with CI \textgreater{} 0; for C1-strict, CI \textgreater{} 0; and for C2, a lower CI bound above $-$3. Family 2 used a non-inferiority margin of $-$2 points against the better of P2Rank 2.5.1 and fpocket+PRANK, with Holm correction over six endpoints [31]. Strata (fold novelty by maximum Foldseek TM-score to training, ligand class, pLDDT band and dimer site class) were pre-registered as descriptive. Analyses labelled post hoc were added after the test runs. The pre-registered sensitivity analysis on the common chain subset was run (C1, v2 seed 0: +11.8 [9.7, 13.8] on 6,700 sites). The pre-registered sensitivity analysis with UniProt accessions as bootstrap units was not run.

\subsubsection*{Pre-registration, amendments and test versions}

The pre-registration was frozen at 15:39 CET on day 1 (git time), before any test-set run. Amendment A1 changed the grouping of symmetric dimer copies to a Jaccard index of 0.5 or more, before any method output existed. Amendment A2 defined how ligand atoms are placed for C1-strict. Amendment A3 changed the fpocket version of the reproduction gate to 4.2.3, which matches the LIGYSIS authors' pocket counts. A2 and A3 were made before any BURROW test prediction. Test v1 was the frozen model at feature freeze. Test v2 added the output floor (decision D13). The planned third version (v3, interface sites) and a later per-region floor (v3b) both failed their pre-fixed dev gates, so v2 is final. Twelve models (3 seeds $\times$ 2 rankers $\times$ 2 presets) were frozen per version.

\section*{2. Extended results}

\subsection*{Output length and the closest-pocket rule}

The second property is a rule of the LIGYSIS protocol (main text Fig 4A and 4B). Each site is credited only to the pocket whose centre is closest to the site centre. If a low-ranked pocket lies closer than the ranked hit, the site counts as a miss. Long output lists therefore lose credit. The agents found this on dev data: on cryptic sites in the apo development set (CryptoBench training folds), v1 found a pocket within 12 \AA{} in the top N+2 more often than P2Rank (77.8\% against 72.3\%), yet scored lower under the closest-pocket rule. In 403 of its misses, the pocket closest to the site lay within 12 \AA{} but was ranked beyond N+2; only 5 misses had no pocket within 12 \AA{}. v1 output about 42 pockets per structure on that set, against 4.5 for P2Rank. The v2 output floor shortened the LIGYSIS list from 32.9 to 9.1 pockets per chain. C1 rose from 64.3\% to 72.5\%, while C1-strict (54.4\% to 54.2\%) and P1K (95.7\% to 96.1\%) were unchanged. Coverage fell. Maximum recall over the whole list dropped from 99.3\% to 84.2\%. Recall at any rank within N+2 dropped from 80.1\% to 76.9\%, and one-to-one recall from 76.0\% to 71.3\%. The C1 gain from v1 to v2 is therefore an effect of output length, not of better ranking. The baselines were run with their default list lengths, and they might also gain from a tuned floor. C1-strict is robust to the output floor, but not to ligand class: its overall margin of +22.6 points shrinks to +3.1 on non-ion sites. For exhaustive pocket discovery, \texttt{--min-score none} restores the v1 list exactly.

\subsection*{Runtime}

Fig 5 of the main text shows accuracy against runtime, measured in one session on the same 200 LIGYSIS chains for every tool. We ran each tool as one process over the 200 chains and divided the total time, including start-up and file input and output, by 200. BURROW's default preset took 142 ms per structure, P2Rank 2.5.1 (batch mode, one thread) took 209 ms and fpocket 4.2.3 took 605 ms. Default BURROW is therefore about 1.5 times faster than P2Rank per process, which misses the design target of at least 5 times. An optional fast preset took 48 ms (about 4.4 times faster than P2Rank). It was added after the test runs and was not pre-registered. On dev data it cost 0.8 points of top-N+2 recall and 1.7 points of DCA4 recall against the default. On one warm CPU core, including parsing, the fast preset took 31.5 ms (about 6.6 times less than P2Rank's 209 ms, which includes Java start-up) and the default 99.2 ms. Only the fast preset therefore met the design target of at most 70 ms per structure per core; the default took 96.0 ms (median, excluding parsing). Over 10,000 structures, three GPUs processed 411 structures per second with the default preset, against 171 per second for 24 CPU workers. The node was shared with other users' jobs, and repeated sessions differed by up to about 1.3-fold (P2Rank 270 against 209 ms), so these ratios carry an uncertainty of about $\pm$30\%. The Numba CPU backend and the Triton GPU backend give output identical to the reference implementation. The human AlphaFold proteome (20,550 proteins) ran in 129 s on three GPUs in AlphaFold mode (frozen v1 model, before the final speed work).

\subsection*{AlphaFold mode}

In AlphaFold mode, BURROW removes terminal segments with pLDDT below 50 and internal low-confidence segments of 7 or more residues, following AF2BIND [21]. It then makes atoms with pLDDT below 80 partly transparent, linearly down to full transparency at pLDDT 40, so that low-confidence atoms block fewer rays and do not exclude candidate points. The standard ranker is used unchanged. With uniform pLDDT of 90 or more, the output is byte-identical to the standard mode.

On the AlphaFold test set (2,526 human AFDB models, 6,502 LIGYSIS sites), BURROW in AlphaFold mode reached 64.5\% top-N+2 recall, against 51.4\% for P2Rank's AlphaFold configuration (main text Fig 6A). This difference, too, came mostly from ion sites: +32.1 [28.3, 36.1] on ion sites and +1.9 [$-$0.1, 3.9] on non-ion sites (post hoc). Against BURROW on the unmodified model, AlphaFold mode added +1.0 [0.1, 2.0] points of recall and reduced hallucinated pockets among the top N+2 from 2.80\% to 0.12\%. A hallucinated pocket here means a pocket with mean pLDDT below 50 that matches no site. However, a simple filter was equally good. This filter drops pockets with mean pLDDT below 70, as in HOTPocket and in the work of Zillmer and Walther [19, 20]. AlphaFold mode was ahead of it by only +0.4 [$-$0.5, 1.3] points, and the filter had higher precision (+2.4 points P1K). The same equivalence held on development data. The geometric mechanisms remove hallucinated pockets, which the filter also does, and they add no recall beyond the filter.

Precision on AlphaFold models was a weakness. BURROW's P1K on the AlphaFold set was 58.9\% against 73.9\% for P2Rank. The CHEN11-only ablation model (251 P2Rank training structures, no ion labels) reached 78.2\%, which suggests, but does not show, that the low precision comes from the ion-heavy \texttt{all} training. An optional ranker that uses pLDDT descriptors raised P1K to 73.9\%, but it lowered PocketWatch top-N+2 by 3.4 points and P1K by 14.6 points, where no pLDDT is available (descriptive; this recommendation is informed by test data). Two goals of the prompt were not met. PAE matrices could not be used, because the AFDB file endpoint returned HTTP 403. Side-chain tolerance reduced recall on dev and is off by default. It makes side-chain atoms partly transparent; at an opacity factor $\lambda$ = 0.25 on all side chains, recall fell by 10.9 points. The AlphaFold test set is also biased towards holo-like models: every protein has PDB structures, the median local C$\alpha$ RMSD to the PDB representative is 0.38 \AA{}, and only 95 benchmark sites (89 scored) have mean pLDDT below 70.

\subsection*{Dimers}

BURROW takes any number of chains as input and computes candidates, features and ranks on the whole assembly. Because we found no public benchmark of pockets in dimers, the agents built one from 400 PDB assemblies of LIGYSIS proteins (274 homomeric and 126 heteromeric complexes of 2--8 chains). Symmetric copies of a site are merged. For every method, predicted pockets whose residue sets overlap by a Jaccard index of 0.5 or more are merged, by a rule that does not use the site labels. With full-assembly input, BURROW recalled 84.0\% of interface sites in the top N+2, against 87.7\% for P2Rank ($-$3.7 [$-$8.7, 0.9], inconclusive; Fig 6B). On hetero-interface sites the difference was $-$6.1 (66 sites). On single-chain sites within assemblies BURROW was far ahead (+19.0 [13.2, 25.9]), and over all sites it reached 83.9\% against 72.9\%. Only 9 of the 243 interface sites are ion sites, so the interface endpoint is in effect a test on non-ion sites, where BURROW's margins are small elsewhere too. When each chain was given alone (descriptive), BURROW placed a pocket near 69.6\% of interface sites, against 55.4\% for P2Rank (+14.2 [9.0, 20.1]). These values are counted per chain copy (634 rows), so they cannot be compared directly with the assembly values. Counted per site, with a site found if any isolated chain finds it, BURROW lost little relative to assembly input on the 12 \AA{} criterion ($-$1.2) but 11.1 points on the 4 \AA{} atom-distance criterion. A missing partner chain therefore mostly reduces the precision of localisation. A development cycle that targeted interface sites (v3) did not pass its pre-fixed gate (S1 Text). The interface stratum is small (243 sites, 158 bootstrap units), so CIs are about $\pm$5 points.

\subsection*{The agent-run process}

The project ran for about 22 h (from 14:54 on day 1 to 12:35 on day 2), including a 4.5 h pause and the extension granted by the user, and produced more than 130 git commits. The PI agent launched sub-agents for 25 roles during the project and 5 more for the release repository and this manuscript (main text Table 4), and it logged 20 dated decisions (D0--D19). Independent agents caught several errors that the producing agents had missed. The training-set builder's own leakage check passed, but an independent audit found residual leakage, which led to the removal of 186 entries. An exactness test found that a GPU kernel received a threshold in single precision, which changed outputs only for atoms at pLDDT exactly 70. The external report reviewer found that a filter bug had made one descriptive table use all dimer sites instead of interface sites. Reviewer agents also flagged that the v2 floor gained mainly by shortening the output list. Other failures were operational: a broad process-kill command stopped another agent's test job, and timestamps that agents wrote from memory were 20--50 minutes off. An API usage limit paused all agents for 4.5 h. The pause cost only wall-clock time, because long jobs ran as background processes that wrote their state to disk.

\end{document}
